% 量子力学1 量子力学1-11 問題3（059-ring。当時は 量子力学1-7 問題3） 解いた日 2026-11-04
%% kotae
\toi{1}{○}{記述}
\[ \hat{p}\phi=-\ii\hbar\frac{\dd}{\dd x}\phi=p_{n}\phi \]
\[ \therefore\ \frac{\dd}{\dd x}\phi=\frac{\ii p_{n}}{\hbar}\phi,\qquad\phi=\ee^{\ii p_{n}x/\hbar} \]
$\phi(L)=\phi(0)$ より
\[ \ee^{\ii p_{n}L/\hbar}=\ee^{0}\qquad\therefore\ \frac{p_{n}L}{\hbar}=2n\pi \]
\[ p_{n}=\frac{2n\pi\hbar}{L}\quad\nuke{n=0,\pm1,\ldots} \]
\nuke{$\phi'(L)=\phi'(0)$}
\naze{周期的境界条件のもう1つの条件 $\phi'(L)=\phi'(0)$ に触れていない。$\phi'=(\ii p/\hbar)\phi$ なので，$\phi(L)=\phi(0)$ が成り立てば自動的に成り立つ，と一言書く。$n$ が整数 $0,\pm1,\pm2,\ldots$ であることも書く。}
\toi{2}{△}{立式}
\[ \braket{x}{n}=\nuke{\frac{1}{\sqrt{L}}}\exp\Bigl[\ii\frac{2n\pi}{L}x\Bigr] \]
\naze{規格化をしていない。$\phi=A\,\ee^{2\pi\ii nx/L}$ とおくと $\int_{0}^{L}\abs{\phi}^{2}\dd x=\abs{A}^{2}L=1$ で，$x=0$ で正の実数になる位相をとって $A=1/\sqrt{L}$，$\braket{x}{n}=\ee^{2\pi\ii nx/L}/\sqrt{L}$ である。}
\toi{3}{△}{未着手}
\[ \frac{\hat{p}^{2}}{2m}\phi=-\frac{\hbar^{2}}{2m}\frac{\dd^{2}}{\dd x^{2}}\phi=E_{n}\phi \]
\[ \therefore\ -\frac{\hbar^{2}}{2m}\Bigl(\ii\frac{2n\pi}{L}\Bigr)^{2}\phi=E_{n}\phi \]
\[ \therefore\ \frac{4n^{2}\pi^{2}\hbar^{2}}{2mL^{2}}\phi=E_{n}\phi \]
\[ E_{n}=\frac{2n^{2}\pi^{2}\hbar^{2}}{mL^{2}} \]
\nuke{縮退の様子}
\naze{問われている「縮退の様子」が無い。$E_{n}$ は $n^{2}$ だけで決まるので，$n\neq0$ の準位は $n$ と $-n$（右向きと左向き）の2重に縮退し，$n=0$（$E_{0}=0$，定数の関数）だけが縮退しない。}
\[ p_{1}=\frac{2\pi\hbar}{L}=\qty{6.63e-25}{\kilogram\metre\per\second} \]
\begin{align*}
E_{1}&=\frac{2\pi^{2}\hbar^{2}}{mL^{2}}=\qty{2.41e-19}{\joule}\\
&=\qty{1.51}{\electronvolt}
\end{align*}
\toi{4}{△}{未着手}
\[ E_{n}^{\text{井戸}}=\frac{1}{4}E_{n} \]
\nuke{なれない理由}\ \nuke{運動量について}
\naze{残りの2つの問いが無い。運動量の固有関数 $A\,\ee^{\ii px/\hbar}$ は絶対値が $\abs{A}$ で一定なので，$A\neq0$ なら $\phi(0)=0$ を満たせず，井戸のエネルギーの固有関数になれない。井戸の固有状態は運動量の固有状態ではなく，運動量は（大きさも）1つに定まらない。区間の中は $\pm n\pi\hbar/L$ の平面波の重ね合わせの形だが，区間の外で $0$ なので運動量の分布は $k$ の連続関数になる。いえるのは $\expval{p}=0$ と $\expval{p^{2}}=(n\pi\hbar/L)^{2}$ である。}
%% kizuki
\item \mae{第6章 問題1・11-01} 問われた項目の書き落とし（前回は得られる値と確率を答えとして書かなかった）。今回は (3) の縮退，(4) の理由と運動量について言えることが無い。問題文の「求めよ」「述べよ」「答えよ」に印を付け，答案の最後にその数だけ答えが並んでいるかを数える。
\item 規格化された固有関数を求めるときは，$\phi=A\,\ee^{\ii px/\hbar}$ と定数 $A$ を残しておき，最後に $\int_{0}^{L}\abs{\phi}^{2}\dd x=1$ から $A$ を決める。周期的境界条件は $\phi$ と $\phi'$ の2つで，指数関数では $\phi'=(\ii p/\hbar)\phi$ なので2つ目は1つ目から従う。
\item 運動量の固有関数は絶対値が一定なので，壁の条件 $\phi(0)=0$ を満たせない。井戸の $\sin(n\pi x/L)$ を2つの平面波に分けただけで「$\pm n\pi\hbar/L$ が確率 $1/2$ ずつ」とはいえない（区間の外で $0$ なので運動量の分布は連続）。周期的境界条件では $n$ と $-n$ が同じエネルギーで，$n=0$ だけが縮退しない。
